% Chapter 4

\chapter{Back Matter}%
\label{Chapter4}

\section{Appendix}%
\label{sec:app}

There may be more than one appendix. Each appendix should be
thematically consistent. As an example, you could have one appendix for extra
tables and one for pseudocode. However, it directly depends on your topic and
how you present it. Still, the appendix should not be entirely unlinked to the
main text. A small mention like \quo{more data in Appendix X} is sufficient to
justify an appendix.

Every appendix starts with a \verb|\chapter| command. This command behaves
differently in the appendix than in the main text. The numbering uses letters
instead of numbers. Furthermore, the mini \abb{ToC} is not present here.
Generally, the formatting of such an appendix can be done more freely than in
the main text.  For example, one appendix could entirely consist of tables
without any text except for the captions.

\section{Lists of X}%
\label{opt:list}

There are four \quo{List of X} entries in the back matter. The first two are
abbreviations and symbols. More details on both can be found in
\autoref{sec:abb} and \autoref{sec:sym}.

The other two are figures and tables. Both are in some way standard in a
thesis.  However, the real usage of them is small. Do you know someone that
ever looked for a specific table or figure in these lists? So these two lists
are per default not present in the back matter. However, if you still want to
include them in the thesis, use the class options
\clsopt{figlist}\indexdef{figlist} (for figures) and
\clsopt{tablist} (for tables) \indexdef{tablist}.

\section{Bibliography}%
\label{sec:bib}

The bibliography is created from one BibTeX file. This file is given in the
preamble by the command \verb!\addglobalbib!. It can be changed to anything
else. In fact, the bibliography is created with biblatex, which also accepts
other input formats than BibTeX. For an overview look at \citet{web:biblatex}.

The use of \pkg{biblatex} means that the bibliography is not created with
\prog{BibTeX} but with \prog{Biber}. This tool should be contained in any
reasonably modern \LaTeX{} distribution. However, many editors do not use it
as default so you may need to change the configuration there.

The bibliography displays only those entries that are cited at some point in
the thesis (see \autoref{sec:cite}).  With respect to the base on publications
(see \autoref{sub:baseon}), it may be required to have some works in the file
that are not cited anywhere. Note that the bibliography is created for the
main text and the appendix, so by default all citations must be contained in
this single file. You can, however, include other bib files by adding more
\verb!\addglobalbib!  commands to the main \LaTeX{} file.

By default, the first names are shortened to the first letter. This behavior
can be changed with the class command
\clsopt{fullcitename}\indexdef{fullcitename}. When you use this command, the
result is that the full first names are shown in all citations. Note that in
many cases the downloadable \prog{BibTeX} files only contains the shorted
first name.  When some of the citations use shortened and some complete names,
the result will be an inconsistent look. To prevent this, use the command only
when all you BibTeX entries have full first names.

\subsection{BibTex}%
\label{sub:bibtex}

Even though BibLaTeX is more powerful than BibTex, there may be situations
where you want to use BibTex instead, e.\,g.\ if Biber is not available. For
this case, the class option \clsopt{bibtex} can be used.\indexdef{bibtex}
BibTex keeps more or less all features intact. The style of the bibliography
changes a little and the overall look is less pleasant.  The same holds for
the cite commands. Thus, I clearly recommend to use BibLaTeX. If you still
want to use BibTex, you must also adapt the latexmkrc file to properly
configure latexmk for using BibTex.  Be aware that author style and
annotations also does not work with BibTex.

\subsection{Author Style}%
\label{sub:author}

The author style is the way how BibLaTeX shows the author name in the
bibliography. This is set with the class option
\clsopt{bibauthorstyle}\indexdef{bibauthorstyle}.  There are basically two
options: first the given name, then the family name (\cmd{gf}); or first
the family name, then a comma, and then the given name (\cmd{fg}).

More natural is to have the given name shown first. However, this can be
disturbing when the  bibliography is sorted after the family names. Thus, your
brain may have a hard time to figure out where in the sorting you are. A
possible solution is a mixed version of the style. This third possibility is
to have the first author in \cmd{fg} order and all other authors in
\cmd{gf} order (\cmd{fggf}).

You can use any of the three bold strings to both authorstyle class options.
The based on bibliography (\clsopt{baseonauthorstyle}) uses the \cmd{gf}
order by default because the ordering of the papers in this smaller list is
not so important. The main bibliography (\clsopt{bibauthorstyle}) uses the
\clsopt{fggf} order such that most authors can easily be read but the ordering
of the bibliography can also easily be followed.

\subsection{BibLaTeX Annotations}%
\label{sub:annotation}

BibLaTeX has the possibility to add annotations to a bib entry. These
annotations must then be interpreted from the class. This thesis template
interprets three different author annotations. These are the
\cmd{underlinebase}, \cmd{underlineall}, and \cmd{firstauthor}.

\begin{algorithm}
    \caption{An example how to use author annotations.}%
    \label{alg:annotations}
    \begin{code}{latex}
@article{van2000art,
    title={The art of writing a scientific article},
    author={Van der Geer, J and Hanraads, JAJ and Lupton, RA},
    AUTHOR+an = {1=underlinebase,firstauthor;2=firstauthor},
    journal={J. Sci. Commun},
    volume={163},
    number={2},
    pages={51--59},
    year={2000},
}
@misc{web:promotion2,
    author = "{Fakultät für Mathematik und Informatik}",
    AUTHOR+an = {1=underlineall},
    title = "Erste Änderungssatzung zur Promotionsordnung Fakultät für
              Mathematik und Informatik",
    year = "2017",
    howpublished = "\url{http://db.uni-leipzig.de/bekanntmachung/
                          dokudownload.php?dok_id=5033}",
    note = "[Online; accessed 27-August-2018]",
}
    \end{code}
\end{algorithm}


First, it is important to know how author annotations are added. This is done
with an extra field, the \cmd{AUTHOR+an} field. Examples can be seen in
\autoref{alg:annotations}. The content of the field is a list of entries that
address the list of authors. Thus, every item in the list starts with an
author position (starting with 1). Then follows a \quo{=} and the actual
annotations, each separated by a \quo{,}. The list items are then separated with
a \quo{;}.

As an example, this field: \\
\verb|AUTHOR+an = {1=underlinebase,firstauthor;2=firstauthor}|\\
means that the first author is annotated with \cmd{underlinebase} and
\cmd{firstauthor}. The second author is annotated with
\cmd{firstauthor}. Not every position has to be annotated. Thus,
the field \\\verb|AUTHOR+an = {3=underlinebase}|\\ would also be valid.

Now to the meaning of annotations. The \cmd{underlinebase} and the
\cmd{underlineall} annotations do more or less the same thing. They
underline the author name. A convenient way to use this is to underline your
own name in every paper. This makes it easy for the reader to recognize your
work. The difference between these two commands is that \cmd{underlinebase}
only underlines the name in the \quo{based on} bibliography (see
\autoref{sub:baseon}) and \cmd{underlineall} also underlines it in the main
bibliography. You may choose the preferred style according to your taste.
Underlining in the based on bibliography is always a good idea. In the main
bibliography, the underlined references may be a little lost in all the other
ones that are not from you.  However, if you cite many papers that you have
done some work for but the thesis is not based on them, it could help to show
all the work you have done.

The \cmd{firstauthor} annotation is used to mark authors that share the
first authorship of a paper. In the main bibliography, this gives no extra
information that helps finding the work. However, in the based on
bibliography, the situation is different. Maybe you share the first authorship
of a paper, but you are not at the first position. Then you can use this to
make clear that this work is really from you. Thus, the annotation only
effects the based on bibliography. Furthermore, note that this annotation only
makes sense when more than one author is marked. In papers where only one
first author exists, the annotation shall not be used.

\section{Curriculum Scientiae}%
\label{sec:cs}

The Promotionsordung \citep{web:promotion} requires you to hand in a
curriculum scientiae with you work.  This is a curriculum vitae with a focus
on your scientific work. It is not stated how you hand it in but it makes
sense to make it a part of your thesis. To make it an extra file, some extra
commands are provided by the class.

The first one is the \verb!cscategory! environment. This gives you the
environment where you can put \verb!\csentry! in. For every topic (Education,
Working Experience, IT Knowledge and so on), a new \verb!cscategory!
environment is created. Thus, the title of the topic is given as an argument to
the \verb!\begin{cscategory}{title}! command. Internally, the environment
creates a tabular environment such that a \verb!cscategory! must be complete
on one page.  However, this can be composited by breaking one category
manually in parts. The second part should have a \quo{(continued)} (or
\quo{(ctd.)} for short) in the title.

The \verb!\csentry{left}{right}[secondline]! command has as a first argument a
text that is positioned left on the page. This could be a date of an
employment or a subcategory. The second one is then the text that is placed
right on the page.  This is could be a job description or some skills. When
one line is not enough to describe the entry, there is an optional third
argument that creates a second line on the right.

In case you want to create sub-entries for an entry, such as a grade for an
education or a project of an employment, this can be done with the command
\verb|\bull|. This command is used in the second or third argument of the
\verb!\csentry! command. An example of the complete example is given in
\autoref{alg:cs}.

\begin{algorithm}
    \caption{An example how a curriculum scientiae is created.}%
    \label{alg:cs}
    \begin{code}{latex}
\begin{cscategory}{Working Experience}
    \csentry{since 10/2018}{Tutor}[
        Institute for Computer Science, Leipzig University
        \bull{} Lecture: Algorithms and Data Structures
    ]
\end{cscategory}
    \end{code}
\end{algorithm}

\section{Selbstständigkeitserklärung (Statement of Authorship)}
The last thing that is required by the Promotionsordung~\citep{web:promotion}
is that you sign a Statement of Authorship. This is the last page of your
thesis.  It basically says that you have done no plagiarism.

There are two lines. The first one is for the location and date where you
signed it. The second one is for your signature.

% End of chapters/chapter_4.tex
